\begin{tikzpicture}[
  font=\scriptsize,
  kotak/.style={draw=black!45, rounded corners=2pt, text width=2.6cm,
    minimum height=0.8cm, align=center, inner sep=2pt},
  pemicu/.style={kotak, fill=blue!8},
  hasil/.style={kotak, fill=green!10},
  masalah/.style={kotak, text width=4.4cm, minimum height=0.7cm},
  panah/.style={-{Stealth[length=2mm]}, thick}
]
  \node[font=\scriptsize\bfseries] at (0, 1.55) {Pemicu};
  \node[pemicu] (p1) at (0, 0.9) {Jadwal harian pukul 06.00};
  \node[pemicu] (p2) at (0, -0.1) {Berkas hasil laboratorium tiba};
  \node[pemicu] (p3) at (0, -1.1) {Permintaan dari Medikara+};

  % Tiga kotak bertumpuk menandai instans fungsi yang bertambah dan berkurang
  \node[kotak, fill=orange!12, text width=2.4cm, minimum height=1.6cm]
    at (4.5, 0.05) {};
  \node[kotak, fill=orange!12, text width=2.4cm, minimum height=1.6cm]
    at (4.4, -0.05) {};
  \node[kotak, fill=orange!12, text width=2.4cm, minimum height=1.6cm]
    (fungsi) at (4.3, -0.15)
    {\textbf{Fungsi}\\{\tiny\mdseries 0 instans saat sepi,\\ ratusan saat
     ramai;\\ dibayar per eksekusi}};
  \node[font=\scriptsize\bfseries] at (4.4, 1.55) {Platform penyedia};

  \node[font=\scriptsize\bfseries] at (8.8, 1.55) {Tujuan};
  \node[hasil] (h1) at (8.8, 0.9) {Pesan pengingat ke pasien};
  \node[hasil] (h2) at (8.8, -0.1) {Rekam medis dan SATUSEHAT};
  \node[hasil] (h3) at (8.8, -1.1) {Jawaban ke aplikasi};

  \foreach \p in {p1, p2, p3} {
    \draw[panah] (\p.east) -- (fungsi.west);
  }
  \draw[panah] (fungsi.east) -- (h1.west);
  \draw[panah] (fungsi.east) -- (h2.west);
  \draw[panah] (fungsi.east) -- (h3.west);

  \node[masalah, fill=red!22] at (2.15, -2.35)
    {\textit{Cold start}: jeda ketika instans pertama dinyalakan setelah sepi};
  \node[masalah, fill=gray!35] at (6.95, -2.35)
    {Batas waktu per eksekusi dan ketergantungan pada format pemicu penyedia};
\end{tikzpicture}
